\documentclass[11pt]{article}

\usepackage{amsmath}
\usepackage{amsfonts}
\usepackage{amsthm,thmtools}
\usepackage{tikz}
\usepackage{enumerate}


\newcommand{\batsym}{\begin{tikzpicture}[scale=0.05]
	\draw[fill,black]   (-0.25,1.48) .. controls (-0.1,1.5) and (0.1,1.5) .. (0.25,1.48) -- (0.35,1.92) .. controls (0.425,1.8) and (0.41,1.3) .. (0.45,1.2) .. controls (0.6,1.05) and (1.96,1.05) .. (1.98,2.08) -- (5.93,2.08) .. controls (4.2,1.45) and (4,0.3) .. (4.2,-0.28) .. controls (2.4,-0.09) and (0.4,-0.5) .. (0,-2.052) .. controls (-0.4,-0.5) and (-2.4,-0.09) .. (-4.2,-0.28) .. controls (-4,0.3) and (-4.2,1.45) .. (-5.93,2.08) -- (-1.98,2.08) .. controls (-1.96,1.05) and (-0.6,1.05) .. (-0.45,1.2) .. controls (-0.41,1.3) and (-0.425,1.8) .. (-0.35,1.92) -- (-0.25,1.48);
	\end{tikzpicture}}
\newcommand{\QEB}{\begin{flushright}\batsym\end{flushright}}



\def\bold{\textbf}
\theoremstyle{plain}
\newtheorem{problem}{Problem}
\theoremstyle{remark}
\newtheorem*{idea}{Idea}

\renewcommand*{\proofname}{Solution}


\renewcommand{\qedsymbol}{$\batsym$}

\newcommand {\ds}{\displaystyle\sum}
\newcommand {\nnR} {\mathbb{R}_{\geq 0}}

\newcommand {\Mtt}[4] {\begin{bmatrix}#1 & #2\\ #3 & #4\end{bmatrix}}
\newcommand {\Ctt}[2] {\Mtt {#1}{-#2}{#2}{#1}}
\newcommand {\mbR} {\mathbb R}
\newcommand {\mcC} {\mathcal C}
\newcommand {\mbZ} {\mathbb Z}
\newcommand {\mbN} {\mathbb N}
\newcommand {\mbQ} {\mathbb Q}
\newcommand {\RnQ} {\mathbb R \setminus \mathbb Q}
\newcommand {\mcI} {\mathcal I}

\newcommand {\pwii}[6] {\begin{cases} 
		#1 & #2 \\
		#3 & #4 \\
		#5 & #6
\end{cases}}

\newcommand {\pwi}[4] {\begin{cases} 
		#1 & #2 \\
		#3 & #4
\end{cases}}

\newcommand{\eps}{\varepsilon}
\newcommand{\del}{\delta}

\begin{document}
	\title{Math 3A03 - Tutorial 10 Questions - Winter 2019}
	\author{Nikolay Hristov}
	\date{March 25/27, 2019}
	\maketitle
	
        \begin{problem}Let
          \begin{align*}
            f(x) &= \begin{cases}
              x &\text{if $0\leq x\leq 1$}\\
              x-2 &\text{if $1<x \leq 2$}
            \end{cases}\\
            P_n&=\left\{\frac i n\right\}_{i=0}^{2n}= \left\{0,\frac 1 2, \frac 2 n,
                 \ldots ,\frac i n, \ldots,\frac {2n-1} n, 2\right\} \,.
          \end{align*}
          \begin{enumerate}[(a)]
            
          \item Compute $U(f,P_n)$ and $L(f,P_n)$, you may use the fact that $\sum_{i=1}^n i = \frac{n(n+1)}{2}$, then show that given $\eps>0$ there is a $P_N$ such that $U(f,P_N)-L(f,P_N)<\eps$. Conclude that f is integrable on $[0,2]$.
	
\item Modify the above argument slightly, and conclude $f$ is integrable directly from the definition. i.e., What can you say about $\sup L(f,P)$ and $\inf U(f,P)$?

  \end{enumerate}

\end{problem}

\begin{problem}

\begin{enumerate}[(a)]
		\item Let $f$ be continuous on $(a,b)$ and bounded on $[a,b]$. Prove that $f$ is integrable on $[a,b]$.
		\item Let $f$ be piecewise continuous on $[a,b]$ (meaning there are k points $a=x_1<x_2<...<x_j<...<x_k=b$ with $f$ continuous on $(x_j,x_{j+1})$ and $f$ bounded on $[x_j,x_{j+1}]$, which is to say f makes finitely many finite jumps), prove that $f$ is integrable on $[a,b]$.
	\end{enumerate}
\end{problem}

\begin{problem}
	Let \[f(x) = \pwi{x^2}{x\in \mbQ}{0}{x\notin \mbQ}.\] Prove that $f(x)$ is not integrable on $[0,1]$.
\end{problem}


	
\end{document}